\section{GSSH \textrightarrow{} ITIF Crosswalk}
\subsection{Structural Mapping}
\begin{table}[H]
  \centering
  \begin{tabular}{p{0.3\linewidth}p{0.3\linewidth}p{0.28\linewidth}}
    \toprule
    GSSH Layer & ITIF Response & Mechanism \\
    \midrule
    Real Economy Stagnation & T-value Index & Measures transparency deficit \& fiscal space quality \\
    Monetary Distortion & ITIF Core Rules & Neutralizes opacity in credit operations \\
    Institutional Drift & Governance Layer & Introduces non-political transparency standards \\
    Data Overload & CLEAR Timestamp & Restores verifiable records without personal data \\
    Social Fragmentation & UBI-T & Stabilizes households through transparency-indexed envelopes \\
    \bottomrule
  \end{tabular}
  \caption{Alignment of GSSH constraints to ITIF responses}
\end{table}

\subsection{Transition Logic}
\begin{itemize}
  \item GSSH defines the feedback trap; ITIF supplies the non-political exit architecture.
  \item T-value calibration corrects mispricing and aligns incentives with transparency.
  \item CLEAR acts as a non-intrusive verification layer to prevent relapse into opacity cycles.
  \item UBI-T dampens social instability while integrity measures take effect.
\end{itemize}

\subsection{Indicators}
Key metrics: transparency delta (\(\Delta T\)), institutional latency index, recurrence probability of opacity cycles, sovereign risk spreads vs. T-bands, and household stability dispersion.

\subsection{Reference Diagram}
\placeholderfigure{GSSH\_FiatU\_Map.pdf}{Crosswalk from GSSH constraints to ITIF responses}
